% Generated draft: requires mathematical and source review.
\documentclass[11pt]{article}
\usepackage[margin=25mm]{geometry}
\usepackage{fontspec}
\setmainfont{Latin Modern Roman}
\usepackage{amsmath,amssymb,mathtools,unicode-math}
\setmathfont{Latin Modern Math}
\usepackage{enumitem,needspace,xurl,hyperref}
\setlength{\parindent}{0pt}
\setlength{\parskip}{5pt}
\setlength{\emergencystretch}{3em}
\newcommand{\N}{\mathbb{N}}
\newcommand{\Z}{\mathbb{Z}}
\newcommand{\Q}{\mathbb{Q}}
\newcommand{\R}{\mathbb{R}}
\newcommand{\C}{\mathbb{C}}
\newcommand{\F}{\mathbb{F}}
\newcommand{\A}{\mathbb{A}}
\newcommand{\PP}{\mathbb{P}}
\newcommand{\E}{\mathbb{E}}
\newcommand{\eps}{\varepsilon}
\newcommand{\dd}{\,\mathrm{d}}
\begin{document}
\section*{UnsolvedMath ID: 701356}
\textbf{Source problem code:} catalogue entry
\subsection{Definitions and mathematical statement}
All groups in this statement are finite. A formation is a nonempty class \(\mathfrak F\) of groups, closed under isomorphisms, homomorphic images, and finite subdirect products. It is hereditary if every subgroup of every member of \(\mathfrak F\) belongs to \(\mathfrak F\). Thus non-heredity means
\[
\exists A\in\mathfrak F\;\exists B\leq A:\quad B\notin\mathfrak F.
\]
For an explicit local-formation convention, let \(f\) assign to each prime \(p\) a formation or the empty class. Define \(\operatorname{LF}(f)\) to consist of those groups \(G\) for which
\[
G/C_G(H/K)\in f(p)
\]
for every chief factor \(H/K\) of \(G\) and every prime \(p\mid |H/K|\). Here a chief factor has \(K,H\trianglelefteq G\), \(K<H\), with no normal subgroup of \(G\) strictly between them, and
\[
C_G(H/K)=\{g\in G: ghg^{-1}h^{-1}\in K\text{ for every }h\in H\}.
\]
A formation is local if it equals \(\operatorname{LF}(f)\) for some such \(f\).

For \(L\leq M\), put
\[
\operatorname{Core}_M(L)=\bigcap_{m\in M}mLm^{-1}.
\]
Under the chain convention used for this draft, \(H\leq G\) is \(\mathfrak F\)-subnormal if \(H=G\), or if there is a chain
\[
H=H_0<H_1<\cdots<H_n=G
\]
such that
\[
H_i/\operatorname{Core}_{H_i}(H_{i-1})\in\mathfrak F
\qquad(1\leq i\leq n).
\]
The supplied record does not define this terminology. In particular, it does not say whether each inclusion in the chain must be maximal. The formulation below uses the displayed convention; agreement with the intended source convention requires review, and a maximal-inclusion requirement must not be silently treated as equivalent for a non-hereditary formation.

Write
\[
\mathcal S_{\mathfrak F}(G)=\{H\leq G:H\text{ is }\mathfrak F\text{-subnormal in }G\}.
\]
The subgroup lattice has meet \(U\cap V\) and join \(\langle U,V\rangle\), the subgroup generated by \(U\cup V\). Subject to the terminology qualification above, the question is whether there exists a local formation \(\mathfrak F\) such that
\[
\begin{gathered}
\exists A\in\mathfrak F\;\exists B\leq A:\ B\notin\mathfrak F,\\
\text{and for every finite group }G\text{ and all }U,V\in\mathcal S_{\mathfrak F}(G),\\
U\cap V\in\mathcal S_{\mathfrak F}(G)
\quad\text{and}\quad
\langle U,V\rangle\in\mathcal S_{\mathfrak F}(G).
\end{gathered}
\]
No solubility restriction on the ambient groups is supplied. The record additionally states that a negative answer would solve Kourovka Problem 9.75, because all hereditary local formations with this lattice property are already known. This is supplied contextual information, not an independently verified classification or implication; the statement of Problem 9.75 and the classification references are not supplied.
\subsection{Short English statement}
Does there exist a local formation of finite groups that is not closed under taking subgroups, but whose \(\mathfrak F\)-subnormal subgroups are closed under intersection and generated join in every finite group? The precise convention for \(\mathfrak F\)-subnormality requires source review. The supplied record labels the problem open.
\subsection{Sources}
\begin{itemize}
\item UnsolvedMath local registry: supplied complete source record, permanent ID \(701356\), exact source code \texttt{catalogue entry}, title \emph{Kourovka Notebook Problem 17.54}. Upstream revision: \texttt{d371ee917ac9647fbb06cd6fc529deec6f3dda82}. The record was updated on 2026-10-16 and carries the supplied status \texttt{open}; this timestamp is not a literature-verification date.
\item S. F. Kamornikov, proposer, 2010: Problem 17.54 in \emph{The Kourovka Notebook (Unsolved Problems in Group Theory)}, 17th issue (2010), as identified in the supplied 21st-issue edition (2026), arXiv identifier \texttt{1401.0300v48}. The supplied collection description names E. I. Khukhro and V. D. Mazurov as editors. This bibliographic identification is supplied attribution and has not been independently checked.
\item The supplied entry states that a negative answer would solve Problem 9.75 and that the hereditary local formations having the stated property are already known. Neither the text of Problem 9.75 nor references for that classification are included in the record.
\end{itemize}
\end{document}
